$E_0=10(DG_0)^2$.

\paragraph{Modularity and a uniform conductor bound.}
Choose one number field $C$ containing all $a_\sigma$, the values of
$\beta$, and a splitting field of $B$. For each coefficient prime
$\ell$ choose an embedding of $C$ into $\overline{\mathbb Q}_\ell$
and extend scalars on $V_\ell(A)$ to a finite extension containing
its image. With the constant-coefficient convention above, define
\[
 \rho_\ell(\sigma)
   =\beta(\sigma)^{-1}a_\sigma\mu_\sigma\circ\sigma.
\]
The cocycle identities make this a representation of $G_\Q$.
It commutes with $B$, so matrix units give
$\rho_\ell\simeq r_\ell\oplus r_\ell$ with $\dim r_\ell=2$.
On $G_{L'}$ it is the actual Tate action of $A$.
Faltings' semisimplicity and endomorphism theorems
\cite{faltings1983} show that $r_\ell|_{G_{L'}}$ is absolutely
irreducible: after splitting $B$, its commutant is the scalars.
Lemma~\ref{lem:de-rham-descent} below shows that the representation
is de Rham at $\ell$, by proper-smooth comparison and finite descent
\cite{scholze2013,scholze-erratum}, and that its
two Hodge--Tate weights are $0,-1$ in the Tate convention.
Berger's $p$-adic monodromy theorem \cite{berger2002} makes this
de Rham representation potentially semistable, as required by
Pan's theorem.
It is odd. To see this without an assumption on a residual
representation, compare with Betti homology at a complex
embedding. Take the map on rational Betti homology induced by
antiholomorphic conjugation followed by the holomorphic
quasi-isogeny $\mu_c$, and extend it $\mathbb C$-linearly.
It exchanges the two Hodge summands. Each
summand is a simple module for $B\otimes\mathbb C=M_2(\mathbb C)$.
Consequently its action on the two-dimensional multiplicity
space is not scalar; since $r_\ell(c)^2=1$, its eigenvalues
are $1,-1$.
It is unramified outside $T\cup\{\ell\}$, since there $L$ is
unramified, $A$ has good reduction, and $\beta$ is unramified.

Use the Tate convention in which the cyclotomic character has
Hodge--Tate weight $-1$. Apply \cite[Theorem 1.0.4]{pan2022} at
$\ell=5$; there is no residual-image proviso at this prime.
It gives a cuspidal eigenform up to character twist. The twisting
character is a de Rham subrepresentation of the Hom representation
and is unramified outside finitely many primes. Local algebraicity
\cite[Appendix B, pp.~32--33]{conrad-local-characters} and class field
theory over $\mathbb Q$ make it a finite character times an integral
power of the cyclotomic character: removing its integral local
Hodge--Tate weight leaves finite inertia at five and at every other
prime, hence a finite ray-class character. Comparing $\{0,-1\}$
with the modular weights $\{0,1-k\}$ after that power forces
$k=2$ and power zero. Absorb the finite Dirichlet character into
the eigenform and take its primitive newform $f$. Li's Fricke formula
\cite[p.~296, before Lemma 6]{li-newforms} makes the span of its
coefficient conjugates Fricke-stable. Shimura's quotient construction
and Euler-factor identity \cite[Theorem 1]{shimura-factors} therefore
show that $r_5$ is a constituent of the Tate module of the modular abelian variety
$A_f$, or equivalently $r_5^\vee$ is its cohomological constituent.
